\section{Introduction}
\subsection{REXUS Experiment, Cryovolcanoes}
\subsection{Outline}

\section{CFD}
Steps in develop a CFD simulation:
\begin{itemize}
    \item select the mathematical model, defining the approximation to reality
    \item discretisation of the space and equations
    \item analysis of numerical scheme and stability and accuracy
    \item solve the resulting equations and optimize the convergence rate
    \item post processing the result
\end{itemize}~\cite{wilcoxTurbulenceModelingCFD2006}
\subsection{Balance Equations}
Consider a material volume element $\mathcal{V}(t)$ (where $V(t) \in \Omega \in \mathbb{R}^3$) is a volume of fluid with a velocity field $\mathbf{u}(t, \mathbf{x})$, where $\mathbf{u} \in C^{2}$ parallel 
to the sufficiently smooth boundary surface $\partial\mathcal{V}$. The Reynolds' transport theorem
states for a differentiable scalar field $\phi(t, \mathbf{x})$ associated with the fluid 
\begin{equation}\label{theorem_reynolds_tt}
    \frac{d}{dt} \int_{\mathcal{V}(t)} \phi\,d\mathcal{V} =  \int_{\mathcal{V}(t)}(\frac{\partial \phi}{\partial t} + \nabla \cdot (\phi \mathbf{u}))\,d\mathcal{V} 
\end{equation}
Why useful? Recasted weird integral differentiation and helps study the evolution of physical variables in a fluid flow.
The basic equations of fluid dynamics, the Navier Stokes equations, for the case of an isothermal flow represent the mass and linear momentum conservation laws. 
Describing the motion of a fluid using the
\begin{itemize}
    \item mass density field, $\rho(t, \mathbf{x})$,
    \item velocity field, $\mathbf{u}(t, \mathbf{x})$,
    \item and pressure, $P(t, \mathbf{x})$ 
\end{itemize}
which are functions of the time, $t \in \mathbb{R}$ and the spatial position $\mathbf{x} \in \Omega$.
The law of mass conservation states that the rate of change of mass in an arbitrary material volume equals the rate of mass production in the volume~\cite{TOR?BOOK?}. In a closed system with no phase change
(thus, negating the mass production term), the application of~\ref{theorem_reynolds_tt} to the time rate of change of the total amount of fluid in a material 
volume $\mathcal{V}(t)$ results in the continuity equation.
\begin{equation*}
        \frac{d}{dt} \int_{\mathcal{V}(t)}\rho (t, \mathbf{x})\,d\mathcal{V}= \int_{\mathcal{V}(t)}(\frac{\partial \rho}{\partial t} + \nabla \cdot (\rho \mathbf{u})(t, \mathbf{x}))\,d\mathcal{V} = 0
\end{equation*}
\begin{equation}\label{theorem_cont_eqn}
    \frac{\partial}{\partial t}\rho (t, \mathbf{x}) + \nabla \cdot (\rho (t, \mathbf{x}) \mathbf{u}(t, \mathbf{x})) = 0\,\forall\,t \in (0, T],\,x \in \Omega
\end{equation}
Similarly, considering the linear momentum in the arbitrary material volume $\mathcal{V}(t)$ which can be formulated as 
\begin{equation*}
        \int_{\mathcal{V}(t)} \rho(t, \mathbf{x}) \mathbf{u}(t, \mathbf{x})\,d \mathcal{V},
\end{equation*}
the rate of change of linear momentum in the material volume can be recasted using~\ref{theorem_reynolds_tt} as
\begin{equation}\label{eqn_linearmomentum}
    \frac{d}{dt} \int_{\mathcal{V}(t)} \rho \mathbf{u}(t, \mathbf{x})\,d \mathcal{V} = \int_{\mathcal{V}(t)}(
        \frac{\partial}{\partial t} \rho \mathbf{u}(t, \mathbf{x}) + \nabla \cdot (\rho \mathbf{u}(t, \mathbf{x}) \cross \mathbf{u}^{T}(t, \mathbf{x}))
        )\,d \mathcal{V}.
\end{equation}
Substituting~\ref{eqn_linearmomentum} into Cauchy's momentum equation that describes the momentum transport in a continuum in inertial frames, we get the law of
conservation of linear momentum
\begin{equation}\label{theorem_momentum_eqn}
    \frac{\partial}{\partial t}(\rho(t,x) \mathbf{u}(t, \mathbf{x})) + \nabla \cdot (\rho(t,x) \mathbf{u}(t, \mathbf{x})\mathbf{u}(t, \mathbf{x})) =
    \nabla \mathbf{\mathbb{\sigma}(t, \mathbf{x})} + \rho(t, \mathbf{x}) \mathbf{b}(t, \mathbf{x})
\end{equation}
$\,\forall\,t \in (0, T],\,x \in \Omega$, where $\mathbf{b}$ are body forces and $\mathbf{\mathbb{T}}$ is the Cauchy stress tensor.
The convection term here is nonlinear even for incompressible flows and is responsible for the appearance of turbulence~\cite{hirschNumericalComputationInternal2007}

When the first law of thermodynamics is applied to the material volume, the conservation of energy can be written as
\begin{equation}\label{theorem_energy_eqn}
    \frac{\partial}{\partial t}(\rho(t,x)E) + \nabla \cdot (\rho(t, \mathbf{x}) E \mathbf{u}(t, \mathbf{x}))
    = \nabla \cdot (\mathbf{u}(t, \mathbf{x}) \cdot \mathbf{\sigma}(t,\mathbf{x})) + (\rho \mathbf{b}) \cdot \mathbf{u}(t, \mathbf{x})
        -\nabla\cdot\mathbf{q}(t, \mathbf{x}) + S_{E}
\end{equation}
where $E$ is the total energy per unit mass, $\mathbf{q}(t, \mathbf{x})$ is the heat flux, and $S_{E}$ is an energy source per unit volume.
\subsubsection{Non-dimensionalised Navier Stokes Equations}
\subsection{Consitutive Relations}
\begin{itemize}
    \item Eulerian compressible fluid/ideal compressible fluid: $\mathbf{\mathbb{T}} = -p(\rho)\mathbf{\mathbb{I}}$
    \item Eulerian incompressible fluid/ ideal incompressible fluid: $\mathbf{\mathbb{T}} = -p\mathbf{\mathbb{I}}$
    \item incompressible Navier Stokes fluid/incompressible linearly viscous/Newtonian fluid: $\mathbf{\mathbb{T}} = -p\mathbf{\mathbb{I}} + 2\mu\mathbf{D}$
\end{itemize}~\cite{truesdellIntroductionMechanicsFluids2000}

\subsection{FV Methods and Star-CCM+}
Finite volume methods (FVM) are the popular discretisation method used in computational fluid dynamics. FVMs can be described 
as discrete evolution equations for cell averages. There is freedom in relation to representing the numerical solution. [cite ansorge sonar]
The finite volume method uses the integral form of the conservation equation. The solution domain is divided into a finite number of small control volumes (CVs)
by a grid which defines the control volume boundaries but not the computational nodes. Discrete versions of the integral form of the conservation 
equations are applied to each control volume in order to obtain a set of linear algebraic equations.[cite starccm] FVms can be viewed as being part
of the Petrov-Galerkin family of finite element methods with piecewise constant trial functions, enabling the use of 
unstructured grids and leverage the computational advantages of the finite difference and finite element methods.


\subsection{Star-CCM+}
\subsection{Modelling Turbulence}
\subsection{Multiphase CFD}
A gas satisfies the continuum hypothesis if the Knudsen number  is << 1.
\subsection{Fundamentals}
\subsection{Interfaces}
\subsection{Boundary Conditions}
\subsubsection{Volume of Fluid Method}


\subsection{Validation and Verification}
\subsection{Stefan Problem}

\section{Free Surface Flow Problem}
\subsection{Literature Research}
\subsection{Problem Description}
\subsection{Rayleigh-Taylor Instability}
\subsection{Impinging Droplets}
\subsection{Dam Break}
\subsection{}


